\SourceLeaf{178}ständniss scheint uns der Würde der Wissenschaft angemessener, als das eitele
Bemühen, die Lücke, die man nicht ausfüllen kann, durch ein unhaltbares Gewebe
von Scheinbeweisen zu verbergen.

Der Verfasser der erstern Schrift hatte bereits vor 15 Jahren in einer kleinen
Abhandlung: »Tentamen novae parallelarum theoriae notione situs fundatae«
einen ähnlichen Versuch gemacht, indem er alles auf den Begriff von Identität
der Lage zu stützen suchte. Er definirt Parallellinien als solche
gerade Linien, die einerlei Lage haben, und schliesst daraus, dass solche Linien
von jeder dritten geraden Linie nothwendig unter gleichen
Winkeln geschnitten werden müssen, weil diese Winkel nichts anderes
  seien, als das Mass der Verschiedenheit der Lage dieser dritten Linie von den
Lagen der beiden Parallellinien. Diese Beweisart ist in der vorliegenden neuen
Schrift wiederholt, ohne dass wir sagen könnten, dass sie durch die
eingewebten philosophischen Betrachtungen an Stärke gewonnen hätte. Der
Behauptung S. 24: »Notionem situs e geometria adeo non excludi posse, ut potius
notionibus ejus fundamentalibus annumeranda sit, \textls{dudum omnes agnovere
geometrae}« muss in dem Sinne, in welchem der Verf. den Begriff Lage in seinem
Beweise gebraucht, jeder Geometer widersprechen. Wenn wir von des
Verfassers Definition: »Situs est modus, quo plura coëxistunt vel juxta se
existunt in spatio« ausgehen, so ist Lage ein blosser Verhältnissbegriff, und man
kann wohl sagen, dass zwei gerade Linien $A,B$ eine gewisse Lage gegen einander
haben, die mit der gegenseitigen Lage zweier andern $C,D$ einerlei ist. Aber der
Verf. gebraucht das Wort Lage in seinem Beweise als absoluten Begriff, indem er
von Identität der Lage zweier nicht coincidirenden geraden Linien spricht. Diese
Bedeutung ist offenbar so lange leer und ohne Haltung, bis wir wissen, was wir
uns bei einer solchen Identität denken und woran wir dieselbe erkennen
sollen. Soll sie an der Gleichheit der Winkel mit \textls{einer} dritten
geraden Linie erkannt werden, so wissen wir ohne vorangegangenen Beweis noch
nicht, ob eben dieselbe Gleichheit auch bei den Winkeln mit einer vierten geraden
Linie Statt haben werde: soll die Gleichheit der Winkel mit \textls{jeder} andern geraden Linie das Criterium sein, so wissen wir wiederum nicht, ob
gleiche Lage ohne Coincidenz möglich ist. Wir stehen mithin
\textls{nach} des Verf. Beweise noch gerade auf demselben Punkte, wo wir
\textls{vor} demselben standen.

